\documentclass[11pt,letterpaper]{article}
\usepackage[utf8]{inputenc}
\usepackage[spanish]{babel}
\usepackage[margin=2.5cm]{geometry}
\usepackage{amsmath}
\usepackage{amsfonts}
\usepackage{amssymb}

\begin{document}

\begin{center}
    \textbf{\Large Examen de Evaluación: Unidad III (Fluidos)}\\
    \vspace{0.2cm}
    \textbf{Física para Ciencias de la Salud (005-1713)}\\
\end{center}

\vspace{0.5cm}
\noindent \textbf{Nombre y Apellido:} \underline{\hspace{7cm}} \hfill \textbf{C.I.:} \underline{\hspace{4cm}}

\vspace{0.5cm}
\noindent Resuelva los siguientes problemas, justificando analítica y matemáticamente cada uno de sus procedimientos. Exprese sus resultados finales con las unidades correctas del Sistema Internacional (S.I.). Considere la aceleración de la gravedad como $g = 9.8 \text{ m/s}^2$.

\vspace{0.5cm}

\begin{enumerate}

    % Problema 1: Hidrostática aplicada a transfusiones (Despeje de altura con contrapresión)
    \item \textbf{Presión Hidrostática y Venosa:} Un paciente gravemente herido necesita una transfusión de sangre rápida (densidad de la sangre transfusora $\rho = 1060 \text{ kg/m}^3$). Si la presión venosa del paciente en el punto de inserción es de $2500 \text{ Pa}$ (presión manométrica), ¿a qué altura mínima por encima de la vena debe colgarse la bolsa de fluido para que la sangre logre vencer la presión venosa e ingresar al torrente sanguíneo?
    \vspace{0.8cm}

    % Problema 2: Arquímedes inverso (Cálculo de volumen y densidad)
    \item \textbf{Principio de Arquímedes:} En un laboratorio de biomecánica, una muestra de tejido óseo compacto se pesa en el aire arrojando un valor de $4.50 \text{ N}$. Posteriormente, se sumerge completamente en agua destilada ($\rho = 1000 \text{ kg/m}^3$) y su peso aparente medido es de $2.10 \text{ N}$. Utilizando la fuerza de empuje, calcule primero el volumen de la muestra y luego determine la densidad real de dicho tejido óseo.
    \vspace{0.8cm}
    
    % Problema 3: Continuidad + Bernoulli (Estenosis)
    \item \textbf{Hemodinámica (Continuidad y Bernoulli):} La sangre fluye a $0.30 \text{ m/s}$ por una arteria horizontal que posee un radio de $4 \text{ mm}$. Debido a una acumulación de placa aterosclerótica (estenosis), el radio se reduce a la mitad ($2 \text{ mm}$) en un tramo de la arteria. Si la presión en la zona sana es de $14000 \text{ Pa}$, determine la presión en la zona estrechada. (Asuma la densidad de la sangre como $1060 \text{ kg/m}^3$).
    \vspace{0.8cm}
    
    % Problema 4: Ley de Poiseuille (Proporcionalidad y porcentajes)
    \item \textbf{Ley de Poiseuille (Análisis Proporcional):} Debido a la acción de un fármaco vasoconstrictor, el radio de las arteriolas de un paciente disminuye en un $20\%$ respecto a su tamaño original (es decir, el nuevo radio es el $80\%$ del inicial). Suponiendo que la diferencia de presión y la viscosidad de la sangre se mantienen rigurosamente constantes, ¿cuál será el porcentaje exacto de caudal sanguíneo ($Q$) que se mantendrá circulando en comparación con el caudal original?
    \vspace{0.8cm}

    % Problema 5: Ley de Poiseuille (Cálculo completo con despeje de presión)
    \item \textbf{Flujo Viscoso y Resistencia:} Para administrar un medicamento altamente viscoso ($\eta = 2.5 \times 10^{-3} \text{ Pa}\cdot\text{s}$), se emplea una aguja hipodérmica de $0.04 \text{ m}$ de longitud y un diámetro interno de $0.8 \text{ mm}$. Si el tratamiento exige suministrar la medicación manteniendo un caudal rigurosamente constante de $2.0 \times 10^{-7} \text{ m}^3\text{/s}$, ¿qué diferencia de presión ($\Delta P$) debe ser capaz de generar la bomba de infusión mecánica?
    \vspace{0.8cm}

\end{enumerate}

\end{document}